\documentclass[11pt]{article}
\usepackage[T1]{fontenc}
\usepackage[margin=1in]{geometry}
\usepackage{amsmath,amssymb,amsthm}
\usepackage[colorlinks=true,linkcolor=blue,citecolor=blue,urlcolor=blue]{hyperref}
\newtheorem{theorem}{Theorem}[section]
\newtheorem{lemma}[theorem]{Lemma}
\newtheorem{proposition}[theorem]{Proposition}
\newtheorem{definition}[theorem]{Definition}
\newtheorem{assumption}[theorem]{Assumption}
\theoremstyle{remark}
\newtheorem{remark}[theorem]{Remark}
\newcommand{\Tr}{\operatorname{Tr}}
\newcommand{\Erec}{E_{\rm rec}}

\title{Petz Recoverability in AQFT via Conditional Expectations\\
{\large A framework and a conditional exponential recovery bound}\\
{\normalsize Version 2}}
\author{Lluis Eriksson\\ Independent Researcher\\ \texttt{lluiseriksson@gmail.com}}
\date{January 2026 (v1); July 2026 (v2)}

\begin{document}
\maketitle

\begin{abstract}
We formulate an operational notion of recoverability in algebraic quantum
field theory for type III local von Neumann algebras. Fixing a faithful
normal KMS reference state and assuming a state-preserving conditional
expectation, we define the recovery channel and, working in a fixed split
implementation for each separation $r$, we assume (i) exponential decay of
split-implemented conditional mutual information and (ii) a
CMI-to-recovery inequality. Under these explicit bridge assumptions we
obtain a conditional exponential recoverability bound $\Erec(r) \le
g(C_1 e^{-m_{\rm CMI} r})$. \emph{v2 corrects the duality underlying the
construction:} v1 defined the Petz-type channel as the
``Accardi--Cecchini adjoint'' via the pairing $\omega(Z\,\mathcal{R}(X))
= \omega(\varepsilon(Z)X)$ and ``proved'' finite-dimensional consistency
with the standard Petz map; that identity is \emph{false in general} (the
printed proof contains an invalid cyclicity step; numerically the
identity fails at order $10^{-1}$ on random faithful states, holding only
in commuting/product situations). The correct statement, proved and
machine-verified here, is that the standard Petz map is the
\emph{trace-predual} of the generalized (Accardi--Cecchini) conditional
expectation; accordingly, v2 defines the recovery channel in
Schr\"odinger picture as precomposition with the conditional expectation.
This also repairs a type/direction error in v1's recovered-state
definition, whose corrected form $\tilde\omega := \omega|_{AB} \circ
\Psi_r$ (with $\Psi_r$ the normal extension of $\mathrm{id}_A \otimes
\varepsilon$) reproduces the standard Petz reconstruction exactly in
finite dimensions. We further relabel v1's finite-dimensional map as the
generalized conditional expectation (its Takesaki module property fails
generically -- verified), record that for a \emph{true} state-preserving
conditional expectation the recovery is the CE-pullback, and add series
positioning: this note is the AQFT capstone announced by the companions,
its split-implemented CMI is one of three compatible regularizations in
the series, and the numerical program deferred by v1 has since been
executed. The main theorem is unchanged: a conditional framework
statement isolating the missing bridge assumptions.
\end{abstract}

\section{Introduction}
Local algebras in AQFT are typically type III, so reduced density
matrices and partial traces are unavailable in the naive sense. This note
is intentionally conservative: assuming quantitative conditional
independence and a CMI-to-recovery bridge (both explicit), one obtains an
exponential bound on a fidelity-based recovery error for a recovery
channel built from a conditional expectation. The main statement is a
conditional theorem isolating AQFT-correct infrastructure; the bridge
assumptions are the open technical problem.

\paragraph{Series positioning (added in v2).} This is the AQFT companion
announced in 2601.0043/0044 \cite{S0043, S0044} and the sibling of the
Type III capstone 2601.0034 \cite{S0034}: there, recovery is the
universal/twirled modular Petz candidate with an assumed FR-type
inequality in purified distance; here, recovery is built from a
state-preserving conditional expectation in a fixed split implementation.
The split-implemented CMI of Section~\ref{sec:bridge} is the third
regularized CMI of the series -- alongside the embedded-state CMI of
2601.0020 \cite{S0020} and the Araki-difference CMI of 2601.0034 -- and
all three reduce to the standard CMI for natural Type I data. The
``numerical outlook'' deferred in v1 has been executed in the companions:
TFIM pipelines in 2601.0043 v2 and a first $\mathbb{Z}_2$ gauge testbed
in 2601.0044 v2. Upstream: 2512.0060/0101 and 2601.0007
\cite{S0060, S0101, S0007}.

\section{Tripartition geometry and split implementations}
As in v1 (unchanged): Haag--Kastler net, faithful normal KMS state
$\omega_\beta$; tripartition $A$--$B$--$C$ with collar $B$ and separation
$r$; Assumption (fixed split implementation): for each $r$ a normal
$*$-isomorphism $\alpha_r : \mathcal{A}(A\cup B\cup C) \to M_{ABC}(r)
\cong B(\mathcal{H}_A^{(r)}) \bar\otimes B(\mathcal{H}_B^{(r)})
\bar\otimes B(\mathcal{H}_C^{(r)})$ mapping each local algebra into its
factor, with pushed-forward state $\omega^r_\beta$; and generation of the
unions by the corresponding subalgebras.

\section{Finite-dimensional warm-up (corrected in v2)}
\label{sec:warmup}
The CMI-to-recovery prototype is unchanged: by Fawzi--Renner \cite{FR}
there is a recovery channel with $-\log F \le c_{\rm FR}\,
I_\rho(A{:}C|B)$; with the squared-fidelity convention one may take
$c_{\rm FR} = 1$ (this is the importable constant; cf.\ the factor
corrections in 2512.0101 v2 and companions), and the FR map need not be
the Petz map.

\begin{definition}[Generalized conditional expectation; relabeled in v2]
\label{def:gce}
For faithful $\rho_{BC} > 0$ with $\rho_B := \Tr_C\rho_{BC}$, define
$E_{BC\to B} : B(\mathcal{H}_B \otimes \mathcal{H}_C) \to
B(\mathcal{H}_B) \otimes \mathbf{1}_C$ by
\[
E_{BC\to B}(Z) := \rho_B^{-1/2}\, \Tr_C\!\bigl(\rho_{BC}^{1/2} Z
\rho_{BC}^{1/2}\bigr)\, \rho_B^{-1/2} \otimes \mathbf{1}_C .
\]
$E$ is normal, unital, completely positive and $\omega$-preserving
($\omega = \Tr(\rho_{BC}\,\cdot)$), but it is \emph{not} a conditional
expectation in the Takesaki sense: the module property
$E(X_B \otimes \mathbf{1}) = X_B \otimes \mathbf{1}$ and idempotency fail
for generic $\rho_{BC}$ (suite: deviations of order $0.5$ on random
faithful states). It is the \emph{generalized} (Accardi--Cecchini)
conditional expectation \cite{AC}; v1 called it ``the $\omega$-preserving
conditional expectation,'' which we correct.
\end{definition}

\begin{proposition}[Correct duality: Petz is the trace-predual of $E$;
replaces v1's Prop.~3.3]
\label{prop:predual}
With $E_0(Z) := \rho_B^{-1/2}\Tr_C(\rho_{BC}^{1/2}Z\rho_{BC}^{1/2})
\rho_B^{-1/2}$ and the standard Petz map
$\mathcal{R}^{\rm Petz}_{B\to BC}(X_B) = \rho_{BC}^{1/2}
(\rho_B^{-1/2}X_B\rho_B^{-1/2}\otimes\mathbf{1}_C)\rho_{BC}^{1/2}$, one
has, for all $X_B$ and all $Y \in B(\mathcal{H}_B\otimes\mathcal{H}_C)$,
\[
\Tr\bigl[\mathcal{R}^{\rm Petz}_{B\to BC}(X_B)\, Y\bigr] \;=\;
\Tr\bigl[X_B\, E_0(Y)\bigr] :
\]
the Petz map is the adjoint of $E$ with respect to the \emph{trace}
pairing, i.e.\ the predual map $E_*$ on density operators.
\end{proposition}
\begin{proof}
$\Tr[\rho_{BC}^{1/2}(\rho_B^{-1/2}X_B\rho_B^{-1/2}\otimes\mathbf{1})
\rho_{BC}^{1/2} Y] = \Tr[(\rho_B^{-1/2}X_B\rho_B^{-1/2}\otimes\mathbf{1})
\rho_{BC}^{1/2}Y\rho_{BC}^{1/2}] = \Tr_B[\rho_B^{-1/2}X_B\rho_B^{-1/2}
\Tr_C(\rho_{BC}^{1/2}Y\rho_{BC}^{1/2})] = \Tr[X_B E_0(Y)]$. (Verified to
$2\times10^{-15}$ on random states.)
\end{proof}

\begin{remark}[v1's pairing is false in general]
\label{rem:pairingfalse}
v1 defined the ``Petz dual'' by $\omega(Z\,\mathcal{R}(X)) =
\omega(E(Z)X)$ and offered a finite-dimensional proof that the standard
Petz map satisfies it. The proof's cyclicity step (moving a single
$\rho_{BC}^{1/2}$ through the trace) is invalid, and the identity fails
numerically at order $10^{-1}$ on random faithful $\rho_{BC}$ (suite);
the same holds for naive GNS- and KMS-weighted variants. The identity
\emph{does} hold in commuting situations -- e.g.\ for product
$\rho_{BC} = \rho_B \otimes \rho_C$, where the true conditional
expectation exists, its predual is $X_B \mapsto X_B \otimes \rho_C$
(equal to the Petz map there), and v1's pairing is satisfied with
$\mathcal{R}$ the inclusion. All of this is machine-verified.
\end{remark}

\section{Fidelity and recovery error}
As in v1: $F(\varphi,\psi)$ the squared transition probability
(Uhlmann--Araki) for normal states \cite{Araki, AR}, monotone under
normal CP unital maps; $\Erec := -\log F(\omega, \tilde\omega)$.

\section{Recovery channel from a conditional expectation (reformulated
in v2)}
\begin{assumption}[$\omega_\beta$-preserving conditional expectation;
unchanged]
\label{ass:ce}
There exists a normal conditional expectation $\varepsilon_{BC\to B} :
\mathcal{A}(B\cup C) \to \mathcal{A}(B)$ with $\omega_\beta \circ
\varepsilon_{BC\to B} = \omega_\beta|_{\mathcal{A}(B\cup C)}$ (a
nontrivial, Takesaki-type modular-invariance condition \cite{AC}).
\end{assumption}
\begin{definition}[Recovery as CE-pullback; replaces v1's Defs.~5.3/6.1]
\label{def:recovered}
Fix $r$ and let $\Psi_r : \alpha_r(\mathcal{A}(A\cup B\cup C)) \to
\alpha_r(\mathcal{A}(A\cup B))$ be the normal unital completely positive
extension of $\mathrm{id}_A \otimes \alpha_r\varepsilon_{BC\to B}
\alpha_r^{-1}$ provided factor-wise by the split implementation. Define
the recovered state
\[
\tilde\omega^{\,r,ABC}_\beta := \omega^{\,r,AB}_\beta \circ \Psi_r ,
\qquad
\Erec(r) := -\log F\bigl(\omega^r_\beta,\, \tilde\omega^{\,r,ABC}_\beta
\bigr),
\]
computable in the type I split algebra by $*$-isomorphism invariance of
$F$ (v1's Lemma 6.5, unchanged and correct).
\end{definition}
\begin{remark}[What was wrong in v1's Definition 6.1]
v1 built a map $\Phi_r$ on products $XY \mapsto X\,\mathcal{R}^{\rm
Petz}(Y)$, which has type $\alpha_r(\mathcal{A}(A\cup B)) \to
\alpha_r(\mathcal{A}(A\cup B\cup C))$, and then set $\tilde\omega :=
\omega^{AB}\circ\Phi_r$ -- ill-typed, since $\omega^{AB}$ is a functional
on the \emph{domain} of $\Phi_r$, not its range. The corrected
Definition~\ref{def:recovered} composes in the admissible direction, and
in finite dimensions it reproduces the standard Schr\"odinger Petz
reconstruction $(\mathrm{id}_A \otimes E_*)(\rho_{AB})$ exactly (suite:
$3\times10^{-16}$), consistently with Proposition~\ref{prop:predual}.
For a \emph{true} conditional expectation
(Assumption~\ref{ass:ce}) the recovered state is the CE-pullback
$\omega_{AB}\circ(\mathrm{id}\otimes\varepsilon)$, exact precisely on
states for which $\mathcal{A}(A\cup B)$ is sufficient.
\end{remark}

\section{Bridge assumptions and conditional theorem}
\label{sec:bridge}
Unchanged from v1 modulo the corrected $\Erec$: split-implemented CMI
$I^{\omega_\beta, r}(A{:}C|B)$ via density operators in the fixed
implementation (scheme-dependent; an Araki formulation would refine it --
cf.\ the compatibility remark in 2601.0034 v2);
Assumption (CMI decay): $I^{\omega_\beta,r}(A{:}C|B) \le C_1
e^{-m_{\rm CMI} r}$; Assumption (CMI-to-recovery): $\Erec(r) \le
g(I^{\omega_\beta,r}(A{:}C|B))$ with monotone $g$, $g(t)\to0$.
\begin{theorem}[Conditional exponential recoverability; unchanged]
Under the standing assumptions, $\Erec(r) \le g(C_1 e^{-m_{\rm CMI} r})$;
if $g(t) \le ct$ in the relevant range, $\Erec(r) \le (cC_1)
e^{-m_{\rm CMI} r}$.
\end{theorem}

\section{Verification suite (added in v2)}
\texttt{verification/2601-0046/verify\_2601\_0046.py}: identifies the
correct duality pairing (trace-predual exact at $2\times10^{-15}$; v1's
$\omega$-pairing fails at $0.77$; naive GNS/KMS variants fail similarly);
verifies that $E$ of Definition~\ref{def:gce} is unital and
$\omega$-preserving but violates the module property (order $0.5$);
verifies the product-case statements of Remark~\ref{rem:pairingfalse};
verifies that the corrected recovered state equals the standard Petz
reconstruction in finite dimensions; and anchors $c_{\rm FR} = 1$
(squared convention) on random tripartite states with Petz recovery
(40/40).

\section{Discussion and outlook}
As in v1: exponential clustering motivates the CMI-decay assumption; the
nuclearity/modular-nuclearity program (quantitative split, finite-rank
approximants, stability of fidelity under limits) is the route to
\emph{deriving} the bridge assumptions \cite{BDL}. Updated: the numerical
program deferred by v1 is now executed in the companions -- lattice
CMI/Petz pipelines (2601.0043 v2) and a first $\mathbb{Z}_2$
lattice-gauge testbed (2601.0044 v2), whose ladder-level lessons
(separation condition; floor management) transfer to any future
continuum-approximation scheme.

\section*{Note added in v2}
The conditional theorem, the bridge assumptions, and the split
implementation are unchanged. v2 corrects the operational core:
(1) v1's defining pairing for the ``Petz dual'' is false in general
(Remark~\ref{rem:pairingfalse}; invalid cyclicity step in its Prop.~3.3;
refuted numerically at order $10^{-1}$), and the correct duality is the
trace-predual identity of Proposition~\ref{prop:predual};
(2) the recovery channel is accordingly defined as CE-pullback
(Definition~\ref{def:recovered}), which also repairs the type/direction
error in v1's Definition 6.1 and reproduces standard Petz exactly in
finite dimensions;
(3) v1's finite-dimensional map is relabeled as the generalized
(Accardi--Cecchini) conditional expectation, its Takesaki module property
failing generically (verified);
(4) series positioning added (three compatible split-CMI regularizations;
differentiation from 2601.0034; numerical program executed in
2601.0043/0044 v2), plus the verification suite. The bridge assumptions
remain the open technical frontier.

\begin{thebibliography}{12}
\bibitem{Petz} D. Petz, Commun. Math. Phys. \textbf{105} (1986)
123--131.
\bibitem{AC} L. Accardi and C. Cecchini, J. Funct. Anal. \textbf{45}
(1982) 245--273.
\bibitem{Araki} H. Araki, Publ. RIMS Kyoto Univ. \textbf{11} (1976)
809--833.
\bibitem{AR} H. Araki and G. A. Raggio, Lett. Math. Phys. \textbf{6}
(1982) 237--240.
\bibitem{BDL} D. Buchholz, C. D'Antoni, and R. Longo, J. Funct. Anal.
\textbf{88} (1990) 233--250.
\bibitem{FR} O. Fawzi and R. Renner, Commun. Math. Phys. \textbf{340}
(2015) 575--611.
\bibitem{S0034} L. Eriksson, ai.viXra:2601.0034 (2026; v2 2026).
\bibitem{S0043} L. Eriksson, ai.viXra:2601.0043 (2026; v2 2026).
\bibitem{S0044} L. Eriksson, ai.viXra:2601.0044 (2026; v2 2026).
\bibitem{S0020} L. Eriksson, ai.viXra:2601.0020 (2026; v2 2026).
\bibitem{S0060} L. Eriksson, ai.viXra:2512.0060 (2025; v2 2026).
\bibitem{S0101} L. Eriksson, ai.viXra:2512.0101 (2025; v2 2026).
\bibitem{S0007} L. Eriksson, ai.viXra:2601.0007 (2026; v2 2026).
\end{thebibliography}
\end{document}
